\section{Conclusion}
\label{sec:conclusion}

This systematic literature review spans three intersecting fields
that together motivate a distinct research direction grounded in the reviewed
corpus: parameter-efficient fine-tuning
and adapter architectures (Sections~\ref{sec:peft}--\ref{sec:adapter_composition}),
adapter-aware inference systems (Section~\ref{sec:inference_systems}), unified by
mixture-of-experts routing as the selection mechanism (Section~\ref{sec:moe_routing}) and
peer-to-peer and federated learning (Section~\ref{sec:p2p_federated}).
The systematic review followed the PRISMA 2020 guidelines \cite{Page2021}, the
eight-step process of Xiao and Watson~\cite{XiaoWatson2019}, and a single-wave
adaptation of the snowballing procedure of
Wohlin et al.~\cite{Wohlin2014}, yielding a final corpus of 123 included
records (120 distinct papers), assembled from nine seed papers and six
thematic pre-validated groups.

The following subsections summarise the main findings, characterise the review's contribution, and discuss threats to validity and limitations.

\subsection{Summary of Findings}
\label{sec:conclusion:summary}

The review established four key findings:

\begin{enumerate}
  \item \textbf{Adapters are structurally suitable for P2P exchange.} LoRA and
    bottleneck adapters are self-contained, small (of order 2--10~MB for
    $r = 4$--$8$ at GPT-2 scale\footnote{Order-of-magnitude estimate derived
    from architecture parameters, not a measured file size. LoRA on the two
    attention projections most commonly adapted (query and value) adds $2rd$
    trainable parameters per projection, where $r$ is the rank and $d$ the
    hidden size. For GPT-2 medium through XL ($d \approx 1024$--$1600$;
    $24$--$48$ transformer layers) at $r = 4$--$8$ this is $\approx 0.4$--$2.5$
    million parameters, i.e.\ roughly $1$--$10$~MB stored in single precision;
    wider rank, more adapted modules, or lower-precision storage move the figure
    within this band.}), additive to a frozen backbone, and task-specific. These properties
    make them natural transmission units in a P2P network. No reviewed work has operationalised this observation into a complete adapter exchange protocol.

  \item \textbf{Adapter composition is technically feasible without joint training.}
    LoraHub \cite{Huang2023} and LoraRetriever \cite{ZhaoLoraRetriever2024} demonstrate
    that independently trained adapters can be composed effectively using
    weight-combination or retrieval-based selection, without access to any task's
    original training data. Composition quality
    depends on the ability to match adapters to queries based on learned capability
    representations. Model Soups \cite{Wortsman2022} is the underlying weight-averaging
    precedent for this class of methods, though it was demonstrated on full
    fine-tuned vision models averaged from a common initialisation on the same
    task, not on independently trained adapters.

  \item \textbf{Federated PEFT confirms communication efficiency but requires
    centralised aggregation.} The federated PEFT literature (the G6 group in
    Table~\ref{tab:groups}, 30 papers in the final corpus) reports major
    per-round communication savings from exchanging adapter-only updates rather
    than full gradients or weights, for example FedPETuning reports over
    90\% savings at within-5\% accuracy degradation relative to centralised PEFT
    \cite{ZhangFedPETuning2023} and that non-IID performance degradation is
    bounded (Section~\ref{sec:p2p_federated:fed_peft}).
    However, all of the reviewed federated-PEFT systems rely on a central
    aggregation server, leaving fully decentralised adapter aggregation as an
    open problem.

  \item \textbf{P2P distributed ML has the protocol infrastructure but lacks adapter
    semantics.} P2P systems (Petals, Šajina, gossip learning) demonstrate that
    distributed model serving and parameter exchange without a central server are
    technically feasible. None of these systems addresses per-task adapter modules,
    adapter discovery, or multi-task adapter composition.
\end{enumerate}

\subsection{Research Gap Identification and Systematisation}
\label{sec:conclusion:contributions}

The central contribution of this review is the identification and structured
characterisation of a compound research gap at the intersection of three active
fields: parameter-efficient fine-tuning, peer-to-peer and federated learning,
and adapter-aware inference serving. The concept
matrix (Table~\ref{tab:concept_matrix}) shows that none of the 17 systems it
scores, selected as the most developed published exemplars of each of the
five pillars (selection rule in
Section~\ref{sec:synthesis_gap:matrix}), simultaneously satisfies all seven
design dimensions of a fully decentralised adapter-based inference setting. The research gap map (Table~\ref{tab:gap_map})
decomposes this compound gap into eight open research questions across four
quadrants, each demanding a different type of research contribution.

The systematisation of the 123 included records (120 distinct papers) across five thematic clusters, the identification of a theoretical gap concerning adapter reuse quality under non-IID distribution shift, and the identification of the missing
intersection between DHT-based discovery, adapter capability representation,
decentralised fusion, and P2P topology together constitute the analytical
contribution of this review.
Within the reviewed corpus, no prior finding has examined this intersection; the present work establishes the evidence base and the open
questions on which future empirical work can be grounded.

\subsection{Threats to Validity}
\label{sec:conclusion:threats}

Several threats to the validity of this review are acknowledged.

\textit{Construct validity.}
The seven dimensions used in the concept matrix (Table~\ref{tab:concept_matrix}) were derived inductively from recurring architectural properties in the reviewed literature, not from an a priori theoretical framework. A different set of dimensions could yield a different gap profile. To mitigate this risk, the rationale for each dimension's inclusion and exclusion is documented in Section~\ref{sec:synthesis_gap:matrix}, and the binary scoring scheme was chosen to avoid introducing precision that the evidence cannot support.

\textit{Internal validity.}
The review was conducted by a single author, which introduces the risk of selection bias at every screening stage; the author treats this as the principal threat to internal validity and does not claim to have fully eliminated it. No inter-rater reliability statistic is reported for the screening
process, because no independent second human reviewer was available; an
earlier draft of this review reported an automated re-screening
diagnostic as a partial calibration check, but this has been removed,
since a non-independent automated process sharing the reviewer's own
screening pipeline and rubric cannot support a valid claim of inter-rater
agreement. The author treats single-reviewer selection bias, together with the
framing-informed construction of the pre-validated G1--G6 corpora
(Section~\ref{sec:methodology:screening_reviewers}), as the two principal
threats to the internal validity of this review, and does not claim to
have eliminated either. To bound the consequent bias, automated keyword scoring (Layer~2), structured term sets, and AI-assisted consistency checks were employed alongside manual review, every per-record decision is logged in audit files for reproducibility, and over-inclusion is mitigated downstream by the full-text and quality-appraisal stages (Section~\ref{sec:methodology}). The use of automated classifiers may have introduced systematic bias; however, all automated outputs were advisory, and final decisions were made by the author. The year filter ($\geq$2021 for snowballed records) excluded potentially relevant earlier work, though this risk was mitigated by the inclusion of pre-2021 foundational papers through the pre-validated G0--G6 corpora. Additionally, eleven of the 123 included records (8.9\%) are arXiv preprints without a peer-reviewed published venue at the search date (this is a narrower, venue-status count; a broader retrieval-source count of 18 records with \texttt{source = ARXIV} is reported and reconciled in Appendix~\ref{app:search_log:pdf}, \S D.6); these were retained on an external-citation proxy for minimal peer scrutiny (Semantic Scholar, retrieved 2026-04-21; 2--318 citations, median 48), rather than on in-corpus citation linking, to avoid a circular inclusion criterion. No formal threshold has been established in the SLR methodology literature for this criterion, and a different threshold could alter the corpus composition. A related reporting gap concerns the abstract-review stage: of the 552-record abstract-review pool, 165 records were excluded with only a coarse SKIP disposition and no individually coded exclusion reason retained in the pipeline data (Section~\ref{sec:methodology:search}); this is disclosed here as a limitation on the granularity of the audit trail at that specific stage, consistent with PRISMA 2020 Item 16b, rather than a gap in the fact of reporting the exclusion itself. A second, analogous granularity limitation concerns the identification stage of the Undermind-assisted route, to which 105 of the 123 final records are attributed: Undermind returns a relevance-ranked candidate list rather than an exhaustive hit count, and the size of the candidate pool it evaluated per query before the author's manual selection was not separately retained. This route is therefore reported from its six exported ranked lists (377 records before cross-group deduplication, 343 after; per-group funnel in Table~\ref{tab:groups} and Appendix~\ref{sec:appendix:corpus_pipeline}, \S A.1) rather than from a raw candidate pool. Every stage downstream of that exported list is enumerated and traceable record-by-record. By contrast,
 the \textbf{101} records dropped at the eligibility stage (extraction $\rightarrow$
 final list) are fully categorised and enumerated per-record
 (Table~\ref{tab:final_exclusions}, Section~\ref{sec:methodology}),
 satisfying PRISMA 2020 Item 16b's reason requirement for records read in
 full.

\textit{External validity.}
The corpus is exclusively English-language and predominantly NLP-focused. Findings may not generalise to computer vision, multimodal, or speech-processing settings where adapter architectures are also increasingly used. The dominance of ACL, EMNLP, and NeurIPS venues in the corpus reflects the field's publication patterns but may underrepresent relevant systems work published in networking or distributed systems venues not covered by the seed papers.
One scope boundary is worth stating explicitly: systems for decentralised or
communication-efficient training of a \emph{full} model from scratch (for
example \textsc{DiLoCo} or \textsc{SWARM} Parallelism) were not sought, because
this review concerns the exchange and composition of adapters over a frozen,
shared backbone rather than collaborative training of the backbone itself,
for which Petals~\cite{Borzunov2023} is the in-scope peer-to-peer exemplar and
is included. Within the in-scope literature, adjacent recent lines are
represented: heterogeneous-rank federated LoRA
(HetLoRA~\cite{ChoHeteroLoRA2024}, FLoRA~\cite{WangFLoRA2024},
FFA-LoRA~\cite{SunImprovingLoRA2024}) and adapter retrieval beyond
LoraRetriever~\cite{ZhaoLoraRetriever2024}
(PHATGOOSE~\cite{Muqeeth2024}, the LoRA-library approach of Ostapenko et
al.~\cite{Ostapenko2024}) are all in the final corpus.

\textit{Conclusion validity.}
One wave of citation snowballing was performed; saturation was assessed by overlap with the pre-validated corpora but was not formally measured using a stopping criterion such as diminishing marginal inclusion rate. The field is rapidly evolving: papers published after the search date (2026-04-21) may address gaps identified in this review. The gap analysis reflects the state of the literature at the time of the search and should be interpreted accordingly.

As a check on the gap claims in Section~\ref{sec:synthesis_gap}, a structured
disconfirming search was conducted (method, queries, and results in
Section~\ref{sec:methodology:search} and
Appendix~\ref{app:db_queries}, \S~C.4; grey-literature methodology after
Garousi et al.~\cite{Garousi2019}) specifically to seek evidence against this
review's own conclusions: to locate any existing decentralised
adapter-exchange or adapter-marketplace system, regardless of publication type,
that might already satisfy the combination of requirements in
Table~\ref{tab:concept_matrix}. No located artefact provided that combination;
the near-misses are catalogued in Table~\ref{tab:adversarial_search}. Among the
candidates were two adjacent non-academic projects: Bittensor,\footnote{\url{https://www.bittensor.com/}, accessed 2026-08-27; see also~\url{https://arxiv.org/abs/2507.02951}.} a blockchain-based decentralised
marketplace for machine-learning resources (compute, inference, and whole
models) organised into incentivised ``subnets,'' and Platformless
AI,\footnote{\url{https://www.platformlessai.org/}, accessed 2026-08-27.} a
commercial project describing peer-to-peer training and storage of LoRA
adapters with on-chain attestation. Neither meets this review's peer-reviewed
inclusion criteria (Section~\ref{sec:methodology}) and neither is treated as
scientific evidence; they are disclosed to document that the search surfaced
disconfirming candidates and to bound the gap claims accordingly.
Substantively, neither addresses the specific combination examined here:
Bittensor operates at whole-model/subnet granularity without adapter-level
capability discovery or fusion, and Platformless AI is session-scoped, lacks a
published discovery or composition protocol, and has not been peer reviewed or
independently evaluated. Their existence indicates practical, industry-level
demand for the general problem addressed by this review, noted as motivating
context rather than as evidence qualifying the gap claims in
Section~\ref{sec:synthesis_gap}.

\subsection{Limitations and Future Work}
\label{sec:conclusion:limitations}

This work has several limitations. All findings rest on collected
literature evidence and the author's own analytical validation and
identification of a theoretical foundation. The open research directions identified in this review remain at the conceptual and analytical level; no empirical work yet exists to validate the open question of bounding on adapter reuse quality or the feasibility of local fusion weight estimation in a P2P setting.

The concept matrix (Table~\ref{tab:concept_matrix}) scores 17 selected systems
spanning the five pillars rather than every one of the 123 included records; the
central gap claim is therefore quantified over that illustrative cross-section,
not the full corpus. Extending the seven-dimension scoring to all 123 records
(or at least the 48 Tier-1 records) and publishing the full matrix as a
supplementary artefact is identified as future work.

Second, the reviewed corpus is exclusively English-language. The PRISMA 2020 
review process covered only peer-reviewed works written in English. The 
generalisability of adapter reuse to multilingual NLP environments is not 
addressed. This is acknowledged as a scope limitation, though it reflects 
the current state of the field: the overwhelming majority of relevant 
research in this domain is published in English, and the existence of 
significant work in other languages would represent an exception rather 
than the norm.

Third, the incentive and trust dimensions of a P2P adapter exchange remain
unaddressed. How to motivate peers to contribute quality adapters is an 
open problem. How to detect corrupted or unverified adapters remains 
unresolved. Choosing between incentive design and trust mechanisms as the 
primary focus for future investigation is itself a challenge, since both 
dimensions are significant and complex. In the current technological 
landscape, however, the security and trust dimension would take precedence driven by the increasing accessibility of this technology and the 
growing potential for misuse. At the same time, resolving trust would 
ease the incentive problem, since the two are closely coupled and each 
influences the other.

Future work should prioritise four directions emerging directly from the gap
map. First, formal derivation of bounds on adapter reuse quality as a function 
of distribution distance and adapter rank, building on the domain adaptation 
framework of Ben-David et al.~\cite{BenDavid2010}. Second, empirical investigation of whether weight-space capability embeddings provide sufficient task-level signal for similarity-based retrieval across diverse adapter families.
Third, whether gossip-based propagation can sustain adapter metadata consistency under realistic P2P churn and bandwidth constraints — an open question whose answer would determine whether the distributed systems infrastructure reviewed in Section~\ref{sec:p2p_federated} is sufficient for adapter-level coordination. Fourth, whether probe-based local estimation of fusion weights can approximate centralised AdapterFusion quality, and under what conditions on probe set size and adapter diversity the approximation degrades, a question that connects the composition literature (Section~\ref{sec:adapter_composition}) to the federated setting (Section~\ref{sec:p2p_federated}).